\section{Daimler Truck Holding AG (Xetra: DTG)}
\label{sec:daimler}

Kurs \je{\Kurs}{EUR} am \KursDatum\QW{Yahoo Finance, Kursreihe DTG.DE}{quelle:kurse}{20.09.2026},
\Mio{\AktienMio}{} ausstehende Aktien (gewichteter Durchschnitt viertes Quartal
2025)\Q{DaimlerA}{281}, B\"orsenwert \Mio{\Boersenwert}{EUR}.

\subsection{Eigent\"umer}

Zum 31.12.2025 hielt die Mercedes-Benz Group AG \Pz{\MbAnteil} der Stimmrechte, der
Mercedes-Benz Pension Trust e.\,V.\ \Pz{\MbPensionAnteil} und die Kuwait Investment
Authority \Pz{\KiaAnteil}; institutionelle Anleger insgesamt \Pz{\InstitutionelleAnteil},
Privatanleger \Pz{\PrivatanlegerAnteil}\Q{DaimlerA}{4}.

\bk{Die fr\"uhere Muttergesellschaft bleibt mit knapp \Pz{\MbAnteil} der gr\"o{\ss}te
Einzelaktion\"ar, mehr als vier Jahre nach der Abspaltung. Zusammen mit dem eigenen
Pensionstr\"ager der Mercedes-Benz Group kommt der Mercedes-Benz-Komplex auf gut ein
Drittel der Stimmrechte, ohne dass ein Beherrschungsvertrag oder eine
Mehrheitsbeteiligung vorl\"age. Das ist ein struktureller Anker, kein
Kontrollverh\"altnis.}

\subsection{Directors' Dealings}

Die \"uber die offizielle Meldeseite des Unternehmens auffindbaren Meldungen (16.03.2026,
Karl Deppen und Karin Radstr\"om) betreffen den Erwerb von Aktien im Rahmen eines
Phantom-Share-Plans, nicht freie Marktk\"aufe; Preis und Volumen wurden dort nicht
angegeben\QW{Daimler Truck, Managers' Transactions}{quelle:managers}{20.09.2026}.

\bk{Ein planm\"a{\ss}iger Erwerb im Rahmen eines Verg\"utungsprogramms tr\"agt keine
Information \"uber die Einsch\"atzung des Kurses durch die Organe; anders als bei GEA
(sechs freie K\"aufe, kein Verkauf) liefert dieser Befund deshalb keinen Ansatzpunkt f\"ur
das Votum. Als Gap festgehalten statt als Signal gewertet.}

\subsection{Wettbewerbsposition}

\begin{table}[htbp]\centering\small
\caption{Behauptete Vorteile und die Zahl, die sie zeigen w\"urde.}\label{tab:moat}
\begin{tabularx}{\linewidth}{@{}XXl@{}}\toprule
Behauptung (Quelle) & Zahl, die sie zeigen w\"urde & H\"alt?\\\midrule
Marktf\"uhrerschaft bei schweren Lkw in Nordamerika nach Absatz\Q{DaimlerA}{36} &
Marktanteil in Prozent & nicht belegbar, keine Zahl auf der zitierten Seite\\
Diversifikation \"uber vier Fahrzeugsegmente (Trucks North America, Mercedes-Benz Trucks,
Trucks Asia, Daimler Buses) neben dem Segment Financial Services & Umsatzanteil je Segment
& ja: Seite~3 des Berichts\Q{DaimlerA}{3}\\
Eigenes Finanzdienstleistungssegment als Absatzst\"utze & Nettoliquidit\"at der
Industrial Business unabh\"angig davon positiv & ja: \Mio{\NettoliquiditaetInd}{EUR}
2025\Q{DaimlerA}{2}\\
\bottomrule\end{tabularx}\end{table}

\bk{Von drei Behauptungen ist die Marktf\"uhrerschaft nicht mit einer Zahl unterlegt; der
Bericht nennt sie als Anspruch, keinen Marktanteil in Prozent. Belastbar ist die
Nettoliquidit\"at des Industriegesch\"afts: Trotz eines eigenen, kapitalintensiven
Finanzdienstleistungsarms tr\"agt das Kerngesch\"aft selbst keine Nettoschulden.}

\subsection{Mehrjahresreihe}

\reihentabelle

\Mio{\FcfMinJahr}{}: \Mio{\FcfMin}{EUR} war der niedrigste freie Zufluss der Reihe,
\Mio{\FcfMaxJahr}{}: \Mio{\FcfMax}{EUR} der h\"ochste. Umsatz, EBIT, Marge, ROCE, freier
Zufluss und Absatz erreichten ihr H\"ochst 2023 und sind seither in jedem Folgejahr
gefallen.

\bk{Das ist keine unternehmensspezifische Schw\"ache, sondern der Nutzfahrzeugzyklus: Nach
dem Nachholbedarf der Jahre 2022/2023 (aufgestaute Bestellungen aus der
Lieferkettenkrise) ist die Nachfrage nach schweren Lkw insbesondere in Nordamerika
eingebrochen, sichtbar am R\"uckgang des Absatzes auf \Absatz{} Tsd.\ Fahrzeuge 2025 von seinem
Hoch 2023 (Tabelle~\ref{tab:reihe}). Vier Jahre sind zu wenig, um einen vollen Zyklus zu
zeigen; die Reihe beginnt mit dem ersten Berichtsjahr nach der Abspaltung.}

\subsection{Kursverlauf}

\kurstabelle

\kgvtabelle

Das KGV betr\"agt heute \KgvHeute, gegen\"uber einem Median von \KgvMedian{} in den
Jahren \KgvJahre; der heutige Wert liegt an der Spitze der erfassten Jahre
(\Pz{\KgvRang} Rang).

\bk{Hier liegt die auff\"alligste Divergenz dieses Berichts. Der Kurs ist von
\je{\KursStart}{EUR} Ende 2022 auf \je{\Kurs}{EUR} gestiegen, ein Zuwachs von \KursZuwachs. Im selben Zeitraum sind EBIT, Marge, ROCE, freier Zufluss und Absatz seit ihrem
Hoch 2023 gefallen. Der Kurs bewegt sich damit nicht mit den Fundamentaldaten dieses
Berichts, sondern gegen sie: Das KGV hat sich fast verdoppelt (von \KgvMedian{} auf
\KgvHeute), weil der Kurs stieg, w\"ahrend der Gewinn je Aktie von \je{\EpsHoch}{EUR} (2023) auf
\je{\EpsHeute}{EUR} (2025) sank. Eine plausible Erkl\"arung ist, dass der Markt den zyklischen
Tiefpunkt bereits als durchschritten ansieht und auf die Erholung vorausl\"auft; belegt
ist damit nur die Kurs- und die Ergebnisreihe selbst, nicht die Erwartung dahinter.}

\subsection{Verwendung des Zuflusses}

\ausschuettungstabelle

\bk{2022, im ersten Jahr nach der Abspaltung, floss weder eine Dividende noch ein
R\"uckkauf; seit 2023 sch\"uttet Daimler Truck regelm\"a{\ss}ig aus. 2025 \"uberstieg die
Aussch\"uttung (Dividende plus R\"uckkauf) den ausgewiesenen freien Zufluss der Industrial
Business deutlich \"uber die Summe des laufenden Jahres hinaus, was bei konstanter
Fortf\"uhrung nicht durchhaltbar w\"are; dass die Nettoliquidit\"at der Industrial Business
dennoch positiv blieb (\Mio{\NettoliquiditaetInd}{EUR}), zeigt, dass die Differenz aus
dem Bestand finanziert wurde, nicht aus laufendem Zufluss.}

\subsection{Analystenkonsens}

\KonsensAnalysten{} Analysten, Gesamtvotum \enquote{Outperform}, mittleres Kursziel
\je{\KonsensZiel}{EUR}\QW{MarketScreener, Konsens der Analysten zu Daimler Truck Holding
AG}{quelle:konsens}{20.09.2026}, \Pz{\KonsensAufschlag} \"uber dem heutigen Kurs.

\subsection{Votum}
\label{sec:votum}

\schwellentabelle
\umkehrtabellen

\vd{\textbf{Nicht kaufen} zu \je{\Kurs}{EUR}. Schwelle auf zehn Jahre
\je{\SchwelleMEDZehn}{EUR}, bei \Pz{\HuerdeLang} \je{\SchwelleMEDZehnLang}{EUR}; der
heutige Kurs ist das \VielfachesMED-fache des ersten Werts. Halter: das Gesch\"aft ist
zyklisch, nicht angeschlagen; das Votum richtet sich gegen den Preis am zyklischen
Erholungs\-narrativ, nicht gegen das Unternehmen.}

\vd{\emph{Begr\"undung.} Zum heutigen Kurs verlangt die Anlage auf zehn Jahre
\WachstumMEDZehn{} Wachstum des freien Zuflusses p.\,a.\ Der freie Zufluss der Industrial
Business ist \"uber die vier erfassten Jahre gefallen, nicht gewachsen, von seinem Hoch
2023 (\Mio{\FcfHoch}{EUR}) auf \Mio{\FcfInd}{EUR} 2025. Der Kurs setzt damit eine
R\"uckkehr weit \"uber das bisherige H\"ochst der Reihe voraus, nicht nur eine Erholung
auf den zyklischen Spitzenwert.}

\vd{\emph{Was das Votum \"andern w\"urde.} Ein Gesch\"aftsjahr, in dem der freie Zufluss
der Industrial Business deutlich \"uber das 2023er Hoch von \Mio{\FcfHoch}{EUR} steigt, also
ein Beleg daf\"ur, dass der zyklische Aufschwung tats\"achlich begonnen hat, nicht nur
erwartet wird. Ein Kurs nahe \je{\SchwelleMEDZehnLang}{EUR}.}

\vd{\emph{Was dagegen spricht.} Vier Jahre sind eine kurze Reihe f\"ur einen zyklischen
Titel; ein einzelner Tiefpunkt (2025) kann das Basisszenario nach unten verzerren.
Sollte der Nutzfahrzeugzyklus 2026/2027 drehen, w\"are der heutige Kurs eine Vorwegnahme,
kein Fehlpreis. Die Analystenmehrheit (\KonsensAnalysten{} Stimmen, \enquote{Outperform})
teilt diese Erwartung, wenn auch als Sekund\"arquelle ohne eigene Nachpr\"ufung
einzustufen. Ferner beobachtet: Der zyklische Charakter (Lkw-Absatz, keine
Auftragsbestandsgr\"o{\ss}e wie bei GEA) macht einen Vergleich mit einem stetigen
Servicegesch\"aft methodisch unpassend; die Umkehrungen in Tabelle~\ref{tab:umkehr}
gegen die eigene Zyklusspitze zu lesen (statt gegen den Median) w\"are die
mildere Bewertung.}
